\documentclass[11pt,letterpaper]{article}

% =========================================================================
% PAQUETES ESTÁNDAR COMPATIBLES CON OVERLEAF (pdfLaTeX / TeX Live)
% =========================================================================
\usepackage[utf8]{inputenc}
\usepackage[T1]{fontenc}
\usepackage[spanish,es-tabla,es-noquoting]{babel}
\usepackage{amsmath,amssymb,amsfonts}
\usepackage{graphicx}
\usepackage{booktabs}
\usepackage{geometry}
\usepackage{xcolor}
\usepackage{tikz}
\usetikzlibrary{arrows.meta,positioning,shapes.geometric,calc,decorations.pathreplacing,babel}
\usepackage{listings}
\usepackage{tcolorbox}
\tcbuselibrary{skins,breakable}
\usepackage{hyperref}

% Márgenes y diseño de página
\geometry{left=2.5cm, right=2.5cm, top=2.5cm, bottom=2.5cm}

% Paleta de colores institucional y técnica
\definecolor{alunagreen}{RGB}{15, 118, 110}
\definecolor{alunadark}{RGB}{15, 23, 42}
\definecolor{alunacyan}{RGB}{6, 182, 212}
\definecolor{alunapurple}{RGB}{147, 51, 234}
\definecolor{alunagray}{RGB}{241, 245, 249}

% Configuración de enlaces hipertexto
\hypersetup{
    colorlinks=true,
    linkcolor=alunagreen,
    filecolor=alunapurple,      
    urlcolor=alunacyan,
    pdftitle={Manual de Operación y Control Térmico - ALUNA PBR-02},
    pdfauthor={Raúl Baytero, Jhon Cárdenas, Santiago Cavadia, Iván Meza - Universidad del Magdalena}
}

% Configuración de listados de código
\lstset{
    basicstyle=\ttfamily\footnotesize,
    keywordstyle=\color{alunagreen}\bfseries,
    commentstyle=\color{gray},
    stringstyle=\color{alunapurple},
    backgroundcolor=\color{alunagray},
    frame=single,
    breaklines=true,
    captionpos=b,
    numbers=left,
    numberstyle=\tiny\color{gray},
    showstringspaces=false
}

% Estilo de cajas de aviso (tcolorbox)
\newtcolorbox{alertaBox}[1]{
    colback=red!5!white,
    colframe=red!75!black,
    fonttitle=\bfseries,
    title=#1,
    arc=3mm,
    boxrule=1mm
}

\newtcolorbox{notaBox}[1]{
    colback=alunagreen!5!white,
    colframe=alunagreen!85!black,
    fonttitle=\bfseries,
    title=#1,
    arc=3mm,
    boxrule=0.8mm
}

% =========================================================================
% INICIO DEL DOCUMENTO
% =========================================================================
\begin{document}

% -------------------------------------------------------------------------
% PORTADA FORMAL
% -------------------------------------------------------------------------
\begin{titlepage}
    \centering
    \vspace*{1cm}
    
    {\Huge \bfseries \color{alunadark} PROYECTO CAPSTONE I -- CONTROL II}\\[0.4cm]
    {\LARGE \bfseries \color{alunagreen} FOTOBIORREACTOR ALUNA PBR-02}\\[0.2cm]
    {\Large \textit{``Aluna Duna'': Sistema Bioinspirado en el Frailejón de la Sierra Nevada}}\\[1.5cm]
    
    \rule{\linewidth}{0.8mm}\\[0.5cm]
    {\huge \bfseries MANUAL INTEGRAL DE HARDWARE, CONTROL TÉRMICO PI Y MECÁNICA BIOMIMÉTICA}\\[0.3cm]
    {\large Cultivo de la Microalga Marina \textit{Tetraselmis chuii}}\\[0.5cm]
    \rule{\linewidth}{0.8mm}\\[2cm]
    
    \begin{minipage}{0.48\textwidth}
        \begin{flushleft} \large
            \textbf{Integrantes (Equipo G3-ALPHA):}\\
            Raúl Baytero\\
            Jhon Cárdenas\\
            Santiago Cavadia\\
            Iván Meza
        \end{flushleft}
    \end{minipage}
    \hfill
    \begin{minipage}{0.48\textwidth}
        \begin{flushright} \large
            \textbf{Institución:}\\
            Universidad del Magdalena\\
            Facultad de Ingeniería\\
            Semillero ProtopIA\\
            Periodo 2026-II
        \end{flushright}
    \end{minipage}
    
    \vfill
    {\large Santa Marta, Colombia \\ \today}
\end{titlepage}

\tableofcontents
\newpage

% -------------------------------------------------------------------------
% SECCIÓN 1: INTRODUCCIÓN Y REQUISITOS BIOLÓGICOS
% -------------------------------------------------------------------------
\section{Introducción y Marco Biológico: \textit{Tetraselmis chuii}}

El fotobiorreactor \textbf{ALUNA PBR-02} es una plataforma de cultivo automatizada diseñada para maximizar el crecimiento celular y la producción de biomasa de la microalga marina verde \textit{Tetraselmis chuii}. 

\subsection{El Reto de la Temperatura: Realidad Experimental vs. Supuestos}
En la literatura general suele asumirse que las microalgas operan adecuadamente entre $25^\circ\text{C}$ y $28^\circ\text{C}$. No obstante, basándonos en datos empíricos de más de 250 lecturas de laboratorio, se determinó que para \textit{Tetraselmis chuii}:
\begin{itemize}
    \item \textbf{Rango Óptimo Real:} \textbf{$20^\circ\text{C}$ a $22^\circ\text{C}$} (Tolerancia máxima segura: $18^\circ\text{C}$ a $24^\circ\text{C}$).
    \item \textbf{Estrés Biológico Severo:} Temperaturas superiores a $25^\circ\text{C}$ inducen degradación de clorofila y ralentizan la fotosíntesis.
    \item \textbf{Umbral Crítico / Colapso Celular:} Por encima de $35^\circ\text{C}$ ocurre lisis celular irreversible y muerte del cultivo.
\end{itemize}

\subsection{El Desafío Térmico en la Costa Caribe}
En un entorno con temperaturas ambiente de $30^\circ\text{C}$ a $32^\circ\text{C}$ (típicas de Santa Marta), enfriar el reactor a $20-22^\circ\text{C}$ representa un incremento del \textbf{39\% en la demanda de refrigeración} frente a un supuesto de $26^\circ\text{C}$. Mantener este gradiente exige retirar más de $60\,\text{W}$ térmicos netos. La radiación lumínica requerida para la fotosíntesis ($2,000 - 4,000\,\text{lux}$, fotoperiodo de 12 a 18 horas) actúa inevitablemente como una fuente externa de calor que sobrecarga el sistema si no se gestiona activamente.

% -------------------------------------------------------------------------
% SECCIÓN 2: BIOMÍMESIS DEL FRAILEJÓN DE LA SIERRA NEVADA
% -------------------------------------------------------------------------
\section{Biomímesis: El Frailejón de la Sierra Nevada y el Escudo Orbital}

El diseño de ALUNA PBR-02 toma como principio rector la estrategia termorreguladora del frailejón (\textit{Espeletia occulta} subsp. \textit{glossophylla}), endémico de los páramos de la Sierra Nevada de Santa Marta.

\subsection{El Principio de la Naturaleza (Frailejón)}
En los ecosistemas de alta montaña, el frailejón enfrenta una radiación ultravioleta y solar extrema durante el día combinada con vientos secos, mientras que de noche soporta temperaturas de congelación. La planta resuelve este dilema mediante:
\begin{enumerate}
    \item \textbf{Pubescencia reflectiva (Vellosidades blancas):} Sus hojas están cubiertas de una densa capa de vellosidades de color blanco que reflejan la radiación infrarroja y visible no aprovechada, evitando el recalentamiento de los tejidos vivos.
    \item \textbf{Manto aislante de hojas secas (Marcescencia):} Las hojas muertas permanecen unidas al tallo formando un abrigo cilíndrico que amortigua los choques térmicos externos.
\end{enumerate}

\subsection{La Adaptación en ALUNA: ``El Frailejón Inverso''}
En nuestro biorreactor se invierte la necesidad funcional: en lugar de retener calor de noche, requerimos \textbf{desviar la carga radiante diurna} sin ahogar el cultivo en la oscuridad permanente:
\begin{itemize}
    \item \textbf{El Escudo Orbital:} Una pantalla semicilíndrica de tela blanca de alta reflectancia albedo montada sobre un brazo rotatorio impulsado por un servomotor de alto torque (\textbf{MG996R}).
    \item \textbf{Operación Dinámica:}
    \begin{itemize}
        \item \textbf{Modo Fotosíntesis ($0^\circ$):} El escudo se ubica en la cara posterior del reactor; la luz incide directamente en el cultivo para alimentar la fotosíntesis.
        \item \textbf{Modo Escudo Frailejón ($180^\circ$):} Cuando el calor ambiental sobrepasa la capacidad de la celda Peltier, el servo rota la pantalla blanca interponiéndola entre la fuente de luz y la botella, reflejando la radiación térmica y aliviando la carga al lazo de control.
    \end{itemize}
\end{itemize}

% -------------------------------------------------------------------------
% SECCIÓN 3: LISTA MAESTRA DE MATERIALES Y EQUIPOS (BOM)
% -------------------------------------------------------------------------
\section{Lista Maestra de Materiales y Equipos (BOM)}

\begin{table}[h!]
\centering
\caption{Componentes físicos y especificaciones del sistema ALUNA PBR-02.}
\label{tab:bom}
\begin{tabular}{lp{6cm}p{4cm}}
\toprule
\textbf{Componente} & \textbf{Especificación Técnica} & \textbf{Función en el Sistema} \\
\midrule
\textbf{ESP32 DevKitC} & 38 pines, dual-core, WiFi/BLE, 240 MHz & Microcontrolador maestro y adquisición \\
\textbf{2x DS18B20} & Sonda impermeable acero inox, bus 1-Wire & Sensor 1 (Agua) y Sensor 2 (Ambiente) \\
\textbf{Resistencia 4.7 k$\Omega$} & Película de carbón 1/4 W & Resistencia pull-up bus 1-Wire sensor 1 \\
\textbf{TEC1-12706} & Celda Peltier 12V 6A, $P_{\max} \approx 60\,\text{W}$ & Actuador térmico de enfriamiento activo \\
\textbf{Placa de Cobre} & Espesor 1--2 mm, alta conductividad térmica & Interfaz de contacto frío sobre la botella \\
\textbf{Módulo D4184} & Doble MOSFET canal N en paralelo, 400W & Conmutador de potencia PWM para Peltier \\
\textbf{Disipador + Fan} & Torre con base de cobre y ventilador 12V & Evacuación forzada del calor de cara caliente \\
\textbf{Servo MG996R} & Piñonería metálica, torque $11\,\text{kg}\cdot\text{cm}$, 5--6V & Actuador mecánico del Escudo Orbital \\
\textbf{Tela Blanca} & Textil sintético reflectante (poliéster/algodón) & Pantalla de albedo para el escudo \\
\textbf{Fuente 30V / 10A} & Regulada conmutada (ajustada a 12V DC) & Alimentación general del sistema \\
\textbf{Reductor Buck 5V} & Módulo LM2596 o fuente 5V 2A independiente & Alimentación dedicada para servo MG996R \\
\textbf{Pasta Térmica} & Compuesto siliconado con óxido de zinc & Acoplamiento térmico en ambas caras Peltier \\
\bottomrule
\end{tabular}
\end{table}

% -------------------------------------------------------------------------
% SECCIÓN 4: ESQUEMA ELÉCTRICO Y CONEXIONES DE HARDWARE
% -------------------------------------------------------------------------
\section{Arquitectura Eléctrica y Tabla Maestra de Conexiones}

\begin{alertaBox}{Regla de Oro: Tierra Común (GND Reference)}
El pin \textbf{GND} de la ESP32 \underline{debe} estar físicamente unido al pin \textbf{GND} de señal del módulo MOSFET y al \textbf{GND} de la fuente del servomotor. Si las tierras están flotando, las señales PWM no tendrán referencia, provocando oscilaciones destructivas o reseteos continuos.
\end{alertaBox}

\subsection{Tabla de Asignación Pin a Pin}

\begin{table}[h!]
\centering
\caption{Conexiones eléctricas en la placa ESP32 (38 pines).}
\label{tab:conexiones}
\begin{tabular}{llll}
\toprule
\textbf{Señal / Función} & \textbf{Pin Físico ESP32} & \textbf{Ubicación en Placa} & \textbf{Detalle Eléctrico} \\
\midrule
\textbf{Sensor 1 (Agua) - Datos} & \texttt{D15} & Fila Inferior, Pin 3 & Bus 1-Wire con Pull-up $4.7\,\text{k}\Omega$ a 3V3 \\
\textbf{Sensor 1 (Agua) - VCC} & \texttt{3V3} & Fila Inferior, Pin 1 & Alimentación 3.3V lógica \\
\textbf{Sensor 1 (Agua) - GND} & \texttt{GND} & Fila Inferior, Pin 2 & Masa de referencia \\
\midrule
\textbf{Sensor 2 (Aire) - Datos} & \texttt{D13} & Fila Superior, Pin 3 & Bus 1-Wire (Pull-up interno activado) \\
\textbf{Sensor 2 (Aire) - VCC} & \texttt{3V3} & Fila Inferior, Pin 1 & Compartido con Sensor 1 \\
\textbf{Sensor 2 (Aire) - GND} & \texttt{GND} & Fila Inferior, Pin 2 & Compartido con Sensor 1 \\
\midrule
\textbf{MOSFET Peltier - PWM} & \texttt{D25} & Fila Superior, Pin 8 & Salida PWM (1 kHz, 8 bits: 0 a 255) \\
\textbf{MOSFET Peltier - GND} & \texttt{GND} & Fila Superior, Pin 2 & Tierra de control \\
\midrule
\textbf{Servo MG996R - Señal} & \texttt{D27} & Fila Superior, Pin 6 & Salida PWM (50 Hz, pulsos $500 - 2500\,\mu\text{s}$) \\
\textbf{Servo MG996R - GND} & \texttt{GND} & Fila Superior, Pin 2 & Retorno común a fuente 5V del servo \\
\bottomrule
\end{tabular}
\end{table}

\begin{notaBox}{Alimentación del Servomotor MG996R}
El cable rojo del servo \textbf{NUNCA} debe conectarse al pin 3V3 ni al pin VIN de la ESP32. Debe recibir $5.0\,\text{V}$ a $6.0\,\text{V}$ directamente desde una salida regulada de al menos $2.5\,\text{A}$ (como un módulo reductor DC-DC LM2596 conectado a la fuente de 12V).
\end{notaBox}

\subsection{Esquema Funcional de Flujo Eléctrico}
\begin{center}
\begin{tikzpicture}[node distance=1.8cm, auto, >=Latex]
    \node [draw, rectangle, fill=alunagray, rounded corners, minimum width=3cm, minimum height=1cm] (fuente) {\textbf{Fuente 12V / 10A}};
    \node [draw, rectangle, fill=alunagreen!20, rounded corners, minimum width=2.5cm, right=2cm of fuente] (mosfet) {\textbf{MOSFET D4184}};
    \node [draw, rectangle, fill=cyan!20, rounded corners, minimum width=2.5cm, right=2cm of mosfet] (peltier) {\textbf{Peltier TEC1-12706}};
    \node [draw, rectangle, fill=alunapurple!20, rounded corners, minimum width=3cm, below=2cm of fuente] (esp32) {\textbf{ESP32 Master}};
    \node [draw, rectangle, fill=orange!20, rounded corners, minimum width=2.5cm, right=2cm of esp32] (servo) {\textbf{Servo MG996R}};
    \node [draw, rectangle, fill=teal!20, rounded corners, minimum width=2.5cm, below=1.5cm of esp32] (sensores) {\textbf{Sensores DS18B20}};

    \draw [->, ultra thick, color=red] (fuente) -- node[above]{\small 12V Potencia} (mosfet);
    \draw [->, ultra thick, color=cyan] (mosfet) -- node[above]{\small PWM Conmutado} (peltier);
    \draw [->, thick, color=blue] (esp32) -- node[left]{\small D25 (PWM)} (mosfet);
    \draw [->, thick, color=purple] (esp32) -- node[above]{\small D27 (50Hz)} (servo);
    \draw [<->, thick, color=teal] (esp32) -- node[right]{\small D15 / D13 (1-Wire)} (sensores);
    \draw [dashed, thick] (fuente) -- node[left]{\small GND Común} (esp32);
\end{tikzpicture}
\end{center}

% -------------------------------------------------------------------------
% SECCIÓN 5: DISEÑO MECÁNICO DEL ESCUDO ORBITAL
% -------------------------------------------------------------------------
\section{Diseño Mecánico y Construcción del Escudo Orbital}

El mecanismo de escudo orbital aprovecha los $180^\circ$ de rotación nativa del servomotor MG996R para rodear la botella cilíndrica del reactor.

\subsection{Materiales de Construcción}
\begin{itemize}
    \item Una botella cilíndrica transparente de 1 a 2 litros (cuerpo del reactor).
    \item Una lámina de plástico semirrígido o el corte longitudinal de una botella plástica de diámetro ligeramente mayor al reactor (radio $+1.5\,\text{cm}$).
    \item Un retazo de tela blanca de alta densidad reflectiva.
    \item Un brazo de servo (horn) de doble salida en plástico o aluminio.
    \item Pegamento epóxico o silicona termofusible.
\end{itemize}

\subsection{Instrucciones de Ensamblaje Paso a Paso}
\begin{enumerate}
    \item \textbf{Conformación del Semicilindro:} Corta la lámina plástica con una altura igual a la zona líquida de la botella y un arco que cubra exactamente $180^\circ$ de la circunferencia.
    \item \textbf{Revestimiento Reflectivo:} Pega la tela blanca estirada sobre la cara externa del semicilindro plástico. Asegúrate de que no queden pliegues sueltos.
    \item \textbf{Montaje del Eje Superior:} Fija un travesaño rígido en la parte superior del semicilindro. En el centro exacto de este travesaño se acopla el aspa del servomotor.
    \item \textbf{Alineación con el Biorreactor:} El servomotor MG996R se fija rígidamente a un soporte situado sobre la tapa de la botella, de modo que el eje del motor coincida con el eje axial central del cilindro de agua.
    \item \textbf{Calibración de Topes:} 
    \begin{itemize}
        \item Al enviar el comando $C:0$, el escudo debe quedar completamente detrás del reactor (frente a la placa de cobre), dejando el $100\%$ del frente despejado a la luz.
        \item Al enviar el comando $C:180$, el escudo debe girar al frente, interponiéndose entre la lámpara y el cultivo.
    \end{itemize}
\end{enumerate}

% -------------------------------------------------------------------------
% SECCIÓN 6: TEORÍA DE CONTROL Y SINTONIZACIÓN PI
% -------------------------------------------------------------------------
\section{Modelo Matemático y Estrategia de Control PI}

\subsection{Modelo FOPDT (First-Order Plus Dead Time)}
El comportamiento térmico del agua acoplada a la celda Peltier y a la placa de cobre se modela mediante una función de transferencia de primer orden con retardo:
\begin{equation}
    G(s) = \frac{T(s)}{U(s)} = \frac{K}{\tau s + 1} e^{-L s}
\end{equation}
Donde:
\begin{itemize}
    \item $K$: Ganancia estática del sistema ($^\circ\text{C} / \text{unidad PWM}$). Dado que el actuador enfría, $K < 0$.
    \item $\tau$: Constante de inercia térmica del agua en segundos (típicamente entre $300\,\text{s}$ y $900\,\text{s}$ según el volumen).
    \item $L$: Tiempo muerto o retardo de transporte térmico a través de la pared de la botella y la placa de cobre (típicamente $10 - 40\,\text{s}$).
\end{itemize}

\subsection{Sintonización de Cohen-Coon para Control PI}
Debido a la presencia del retardo $L$, las fórmulas de Cohen-Coon ofrecen un ajuste robusto frente a perturbaciones ambientales:
\begin{align}
    K_p &= \frac{\tau}{|K| \cdot L} \left( 0.9 + \frac{L}{12\tau} \right) \\
    T_i &= L \left( \frac{30\tau + 3L}{9\tau + 20L} \right) \\
    K_i &= \frac{K_p}{T_i}
\end{align}

\subsection{Control Coordinado por Etapas (Split-Range / Frailejón)}
El algoritmo implementado en el backend ejecuta una estrategia coordinada entre dos actuadores:

\begin{equation}
    \text{Acción}(T_{\text{agua}}) = 
    \begin{cases}
        \text{PWM Modulado por PI}, \quad \theta_{\text{servo}} = 0^\circ & \text{si } T_{\text{agua}} \le 22.0^\circ\text{C} \quad (\text{Zona Normal}) \\
        \text{PWM} = 255 \, (100\%), \quad \theta_{\text{servo}} = f(T) & \text{si } 22.0 < T_{\text{agua}} < 24.0^\circ\text{C} \quad (\text{Alerta}) \\
        \text{PWM} = 255 \, (100\%), \quad \theta_{\text{servo}} = 180^\circ & \text{si } T_{\text{agua}} \ge 24.0^\circ\text{C} \quad (\text{Crítico / Sombra Total})
    \end{cases}
\end{equation}

% -------------------------------------------------------------------------
% SECCIÓN 7: ARQUITECTURA DE SOFTWARE
% -------------------------------------------------------------------------
\section{Arquitectura de Software y Protocolos de Comunicación}

El sistema opera bajo una arquitectura de tres capas desacopladas:

\subsection{Capa 1: Firmware ESP32 (Hardware Abstraction)}
\begin{itemize}
    \item Lee periódicamente los sensores en \texttt{D15} y \texttt{D13} mediante la librería \texttt{DallasTemperature}.
    \item Modula la potencia de la celda en \texttt{D25} usando el periférico \texttt{LEDC}.
    \item Posiciona el servomotor en \texttt{D27} mediante \texttt{ESP32Servo}.
    \item Emite por el puerto USB (\texttt{115200 baud}) una trama CSV cada 2 segundos:
    $$\text{Trama:} \quad \texttt{tiempo,pwm,temp\_agua,temp\_aire,angulo\_servo}$$
\end{itemize}

\subsection{Capa 2: Backend Python (Flask API \& Control Engine)}
\begin{itemize}
    \item Lee el puerto serial de forma asíncrona mediante un hilo (\texttt{threading.Thread}).
    \item Ejecuta la identificación no lineal de curvas con \texttt{scipy.optimize.curve\_fit}.
    \item Calcula en tiempo real el lazo PI digital con integración anti-windup.
    \item Provee una API REST con soporte de CORS para el consumo de telemetría y comandos.
\end{itemize}

\subsection{Capa 3: Frontend Astro (Panel de Control Sci-Fi)}
\begin{itemize}
    \item Construido en Astro v4 y estilizado con Tailwind CSS bajo temática de alto contraste.
    \item Graficación vectorial continua con Plotly.js sin recargar la página.
    \item Controles manuales de deslizadores para calibración de corriente (escalones de 25\%, 50\%, 75\%, 100\%) y prueba inmediata del servomotor.
\end{itemize}

% -------------------------------------------------------------------------
% SECCIÓN 8: PROTOCOLO EXPERIMENTAL DE PRUEBAS
% -------------------------------------------------------------------------
\section{Protocolo Experimental y Guía de Puesta en Marcha}

\subsection{Checklist Previo al Encendido}
\begin{enumerate}
    \item Verificar que el tornillo \texttt{VIN+} y \texttt{VIN-} del MOSFET reciban 12V de la fuente.
    \item Verificar que la perilla de corriente de la fuente (C.C.) esté ajustada al menos en $7.0\,\text{A}$.
    \item Verificar que la masa (\texttt{GND}) de la ESP32 esté unida a la masa del MOSFET y a la masa del servo.
    \item Comprobar que el ventilador del disipador gire a máxima velocidad de forma inmediata.
\end{enumerate}

\subsection{Procedimiento de Caracterización Estática de Corriente}
Para conocer la potencia real de enfriamiento:
\begin{enumerate}
    \item Llenar el biorreactor con el volumen de agua de prueba.
    \item Encender la plataforma web en \texttt{http://localhost:4321} y presionar \textbf{CONECTAR}.
    \item Presionar el botón de \textbf{25\%} en el panel de calibración. Esperar 10 minutos hasta que la temperatura del agua se estabilice y anotar el $\Delta T_1$.
    \item Repetir el procedimiento para \textbf{50\%}, \textbf{75\%} y \textbf{100\%}.
    \item Graficar la relación Corriente vs. Temperatura Estacionaria para encontrar el punto de saturación térmica de la celda.
\end{enumerate}

\subsection{Cierre del Lazo y Verificación del PI}
\begin{enumerate}
    \item En la web, presionar \textbf{Calcular Kp / Ki (Cohen-Coon)} para cargar las constantes identificadas.
    \item Presionar \textbf{Activar Control PI Automático (21.0 °C)}.
    \item Acercar una lámpara o foco térmico a la botella: verificar que cuando la temperatura del agua sobrepase los $22.5^\circ\text{C}$, el servomotor empiece a cerrar automáticamente el escudo orbital reflectivo, estabilizando el cultivo en el rango óptimo de \textit{Tetraselmis chuii}.
\end{enumerate}

\end{document}
